\scanpage{361}
\uline{normalisant} le sous-groupe $G\subset\widehat{\mathcal T}^{+}(\pi'_0)$
(\uline{NB} en fait on a $G\subset\mathcal T^{+}(\pi'_0)$ ; \unsure{$G$|\deleted{$G$}} et $\pi_0$
se \unsure{reconstitue|reconstituent} à partir de $\pi'_0$ et de
ce sous-groupe de $\pi'_0$) --- \deleted{\unsure{je crois}}\inserted{on doit}
pouvoir montrer sans trop de mal
que \deleted{\illegible{1}} dans (8), $\widetilde{\widehat{\mathcal T}}^{\unsure{+|\ast}}(\pi_0)$ et $\widetilde{\widehat{\widehat{\mathcal T}}}(\pi_0)$
sont justement les normalisateurs de $G$
dans $\widehat{\mathcal T}^{+}(\pi'_0)$ et dans $\widehat{\widehat{\mathcal T}}(\pi'_0)$ (ce
qui impliquerait bien a), d'ailleurs) ---
mais je ne vois aucune raison plausible
que (b) soit vrai --- ce qui signifierait,
essentiellement, qu'il n'y a pas plus de
``\unsure{transcendance} arithmétique'' définie par le
(``très gros'' !) $\widehat{\widehat{\mathcal T}}(\pi'_0)$ (ou le $\widehat{\mathcal T}(\pi'_0)^{+}$), que
celle définie par le (``bien plus petit'') groupe
$\widetilde{\widehat{\widehat{\mathcal T}}}(\pi_0)=\operatorname{Norm}_{\widehat{\widehat{\mathcal T}}(\unsure{\pi_0|\pi'_0})}(G)$ \dots\ D'ailleurs ces
conditions a) b) ne sont pas impératives
pour que la conjecture fondamentale
\unsure{reliant} les $\Pi_{g,\nu}$ à $\Gamma_0$ soit \unsure{cohérente} ---
\mspage{360}

\scanpage{362}
\mspage{183}
par contre il faudrait absolument avoir
c) pour pouvoir au moins définir
$A(\pi)\unsure{\longrightarrow|{-}}A(\pi')$ et (si $\overline\Sigma,\overline\Sigma'$ se correspondent)
$\Gamma_{\overline\Sigma}\unsure{\longrightarrow|{-}}\Gamma_{\overline\Sigma'}$, dans le cas $\unsure{\pi'_a|\pi'_0}$ \deleted{\unsure{invariant}}\inserted{caractéristique} dans $\pi$
(et dans tous les cas si on admet de plus
a), qui semble \unsure{assez} plausible, et donne
les \uline{injectivités} qu'il faut. Mais on
se demande bien pourquoi un élément
de \deleted{\illegible{1}} $\widehat{\widehat{\mathcal T}}(\pi'_0)$ simplement parce qu'
il \deleted{\unsure{normalise}}\inserted{est dans} le \deleted{$G$ et qu'il \unsure{normalise}}\inserted{normalisateur $R$ de $G$}
\deleted{\unsure{le même}}\inserted{dans} $\widehat{\widehat{\mathcal T}}(\pi'_0)$, \inserted{\unsure{donc aussi $R$,|aussi dans $R$,}} et qu'il normalise
$R\cap\widehat{\mathcal T}^{+}(\pi'_0)$, devrait normaliser aussi
$\widehat{\mathcal T}^{+}(\pi'_0)$ lui-$\overline{\mathrm m}$ ! Il est vrai qu'on a
fait des hypothèses \unsure{draconiennes} au départ
(\unsure{partant} d'un ss-groupe \uline{car.} $\pi'_0$ de $\pi_0$) qui
doivent bien se \unsure{refléter} par des propriétés
particulières de $G\subset\mathcal T^{+}(\pi'_0)$. Peut-être
finalement l'astuce de se ramener à des
ss-groupes caractéristiques pour pouvoir examiner
\mspage{361}

\scanpage{363}
la situation n'est \unsure{elle} pas si astucieuse.
Pour y voir plus clair, on pourrait déjà
\unsure{essayer} de comprendre le cas où (sans
supposer $\pi'$ caract.) on suppose $\pi'$ d'indice
2 dans $\pi$, donc invariant et $G\simeq\mathbb Z/2$, ou
bien on prend \deleted{\illegible{2}} le noyau de l'hom
canonique $\pi\longrightarrow(\pi_{\mathrm{ab}})_2$ ($\simeq(\mathbb Z/2\mathbb Z)^{2g\unsure{+|\deleted{+}}}$ si $\nu=0$,
$\simeq(\mathbb Z/2\mathbb Z)^{2g+\nu-1}$ si $\nu\geq1$), en prenant p.ex.
$\pi=\widehat\Pi_{0,3}$ --- \unsure{situation} de la courbe de
Fermat $x^2+y^2+z^2=0$ (revêtement octaédral
de $\mathbb P^1_{\overline{\mathbb Q}}\setminus\{0,1,\infty\}$ \dots)

Mais admettons provisoirement les \deleted{résultats}\inserted{hypothèses}
d'injectivité (relativement anodines), \uline{plus} la
condition c), pour \unsure{en déduire} des \illegible{1} ---
d'où, si $\overline\Sigma$ et $\overline\Sigma'$ se correspondent, un hom
\inserted{``injectif''}
\[
(9)\qquad \Gamma_{\overline\Sigma}\hookrightarrow\Gamma_{\overline\Sigma'},
\]
et voyons ce qu'on pourrait en tirer.
en admettant même, provisoirement (pour
voir) que (9) est un isom \deleted{\unsure{ce qui}}
c'est à dire \unsure{toute|dans} la \unsure{force|forme} de la conjecture
principale $\Gamma_0\simeq\Gamma_{g,\nu}$.
\mspage{362}

\scanpage{364}
\mspage{184}

\unsure{Soient|Soit} $\pi$ un groupe profini \inserted{à lacets} de type $(g,\nu)$, $\pi'$ de type $(g',\nu')$. Appelons ``correspondance'' entre $\pi$ et $\pi'$ un \deleted{groupe} couple formé d'un $\unsure{G|\pi''}$ profini \inserted{à lacets} \deleted{de type $(g'',\nu'')$}, et d'hom. \inserted{à lacets} $\unsure{G|\pi''}\xrightarrow{p}\pi$, $\unsure{G|\pi''}\xrightarrow{q}\pi'$ qui sont \unsure{donc|des} \unsure{chacun} \unsure{composé|composés} d'un morphisme ``bouchage de trous'', et d'une \deleted{inclusion} isom. sur un \unsure{ss-groupe|ss-gr.} d'indice fini.
\unsure{Soit|Soient} $\unsure{J_p|I''_p}$, $\unsure{J_q|I''_q}\in\unsure{J|I''}=I(\pi'')$ les \inserted{\unsure{ss-|s-}}ensembles de $\unsure{J|I''}$ qui correspondent aux ``trous bouchés par $p$'' resp. par $q$, on suppose \inserted{s'il le faut} $\unsure{J_p|I''_p}\cap\unsure{J_q|I''_q}=\varnothing$ (sinon on pourrait factoriser par un $\unsure{\overline{\mathrm m}}$ quotient $\widetilde{\pi}''$ de $\pi''$). En d'autres termes si $\unsure{\overline{J}_p|\overline{I}''_p}$, $\unsure{\overline{J}_q|\overline{I}''_q}$ sont les complémentaires de $\unsure{J_p|I''_p}$, $\unsure{J_q|I''_q}$ dans $\unsure{J|I''}$ (en \unsure{sorte} qu'on a : $\unsure{\overline{J}_p|\overline{I}''_p}\to I=I(\pi)$, $\unsure{\overline{J}_q|\overline{I}''_q}\to I'=I(\pi')$) on a
\[
\unsure{\overline{J}_p|\overline{I}''_p}\cup\unsure{\overline{J}_q|\overline{I}''_q}=\unsure{J|I''}.
\]

\deleted{\unsure{Utilisant|Utilisons} les ``résultats'' des deux \unsure{paragraphes} \unsure{antérieurs}, on définit alors des applications}

On va supposer maintenant $\pi$ muni d'un $\overline{\Sigma}$ \inserted{$\in A(\pi)$}, $\pi'$ muni d'un $\overline{\Sigma}'\in A(\pi')$, on appellera ``correspondance arithmétique'' entre $\pi$ et $\pi'$ un quadruple $(\unsure{G|\pi''},p,q,\unsure{\overline{\Sigma}_G|\overline{\Sigma}''})$, où $(\unsure{G|\pi''},p,q)$ \unsure{sont|est} comme dessus, et $\unsure{\overline{\Sigma}_G^{a}|\overline{\Sigma}_G}\in A(G)$, tels que l'on ait $\unsure{\overline{\Sigma}_G|\overline{\Sigma}_a}=p^*(\overline{\Sigma})=q^*(\overline{\Sigma}')$. \deleted{\unsure{On supp}} Si $P_{\overline{\Sigma}}$,
\mspage{363}

\scanpage{365}

$P_{\overline{\Sigma}'}$, $P_{\unsure{\overline{\Sigma}_G|\overline{\Sigma}_a}}$ sont resp. les torseurs sur $\Gamma_{\overline{\Sigma}}$, $\Gamma_{\overline{\Sigma}'}$, $\Gamma_{\overline{\Sigma}''}$ qu'on sait, on a \unsure{un|deux} \unsure{diagramme|diagrammes} d'isom. de torseurs
\diagram{\begin{tikzcd}[ampersand replacement=\&,column sep=small,row sep=large]
\& P_{\unsure{\overline{\Sigma}_G|\overline{\Sigma}_a}} \& \\
P_{\overline{\Sigma}} \arrow[ur,"p_P^*","\sim"'] \& \& P_{\overline{\Sigma}'} \arrow[ul,"q_P^*"',"\sim"]
\end{tikzcd}}
\diagram{\begin{tikzcd}[ampersand replacement=\&,column sep=small,row sep=large]
\& \Gamma_{\unsure{\overline{\Sigma}_G|\overline{\Sigma}''_G}} \arrow[dr,leftrightarrow,"q_\Gamma^*","\sim"'] \& \\
\Gamma_{\overline{\Sigma}} \arrow[ur,"\unsure{\widehat{p}_\Gamma|p_\Gamma^*}","\sim"'] \& \& \Gamma_{\overline{\Sigma}''}
\end{tikzcd}},
d'où par composition un \unsure{isom|hom} de
\deleted{$\unsure{P_{\overline{\Sigma}}\to}$}
\[
(P_{\overline{\Sigma}},\Gamma_{\overline{\Sigma}})\xrightarrow{\sim}(P_{\overline{\Sigma}'},\Gamma_{\overline{\Sigma}'})
\]

On appelle ``isom arithmétiques extérieurs'' \inserted{de $(\pi,\overline{\Sigma})$ vers $(\pi',\overline{\Sigma}')$,} tout isomorphisme de torseurs qu'on peut obtenir de cette façon. Deux correspondances sont dites \inserted{\unsure{équivalentes} \deleted{\unsure{du point de}} du point de vue} \deleted{\unsure{arithmétiques|arithmétique} \unsure{extérieures|extérieure}} arithmétique extérieur, si elles définissent le $\unsure{\overline{\mathrm m}|=}$ isom. arithm. extérieur. \unsure{La} composition est définie par composition des \unsure{relations} sur $(P_{\overline{\Sigma}},\Gamma_{\overline{\Sigma}})$ --- on voit donc \illegible{1} provient d'une correspondance. On \illegible{1} donc un groupoïde, dont on vérifie sans peine qu'il est connexe. (\inserted{\unsure{Je dis que}} Les autom de $(\pi,\overline{\Sigma})$ ``sont'' les autom de $(P_{\overline{\Sigma}},\Gamma_{\overline{\Sigma}})$ définis par les éléments de $\Gamma_{\overline{\Sigma}}$. \deleted{C'est \unsure{évident} ou \unsure{pas}} C'est \unsure{facile}\ldots

\margin{On appellera \illegible{4} des ``correspondances extérieures''.}
\mspage{364}

\scanpage{366}
\mspage{185}

Dans ce groupoïde, il y a un syst. transitif d'iso. entre les \deleted{\illegible{1}} $(\widehat{\pi}_0,\overline{\Sigma}_{\pi_0})$ (où $\pi_0$ est un groupe discret à lacets) plus gén. entre les $(\pi,\overline{\Sigma}_\alpha)$, où $\alpha$ \inserted{$\in P(\pi)$} est une prédiscrétification de $\pi$ --- définissant donc un isom. \inserted{arithmétique} extérieur de $\Pi_{g,\nu}$
$\pi$ \ldots\ \unsure{Le|Leur} groupe de \unsure{ces|ses} autom. est noté $\Gamma_0$. La \deleted{catégorie} groupoïde des corres \inserted{virtuelles} arithmétiques extérieures s'identifie donc à la cat. des torseurs sur $\Gamma_0$,
Une prédiscrétification (on pourrait aussi l'appeler une rigidification arithmétique extérieure) \unsure{n'est rien autre chose qu'} un isom arithm. extérieur entre cet objet de référence, et $\pi$.
\mspage{365}

\scanpage{367}
\mspage{186}

\uline{29} Critique de l'approche précédente.

L'approche \unsure{du|des} \S\ \unsure{précédent|précédents} semble \unsure{pourtant|franchement} très brutale. J'ai même des doutes si \unsure{la|l'} \deleted{hypothèse} conjecture sur les propriétés du foncteur "bouchage de trous" est vraie telle quelle. Il est vrai que pour un groupe profini \deleted{\illegible{1}} \inserted{\unsure{à}} \unsure{lacets} $\pi$, si on fixe $i\in I(\pi)$, on a un hom canonique
\[
\tag{1}
\widehat{\widehat{\mathcal T}}(\pi,i)
=\underset{\substack{\text{stabilisateur}\\\text{de }i\text{ dans }\widehat{\widehat{\mathcal T}}(\pi)}}{\text{\uline{$\widehat{\widehat{\mathcal T}}(\pi)_i$}}}
\longrightarrow \widehat{\widehat{\mathcal G}}(\pi'),
\]
mais il est problématique si c'est un isomorphisme --- et même si c'est injectif, ou si c'est surjectif. Mais, \unsure{choisissant} une discrétification $\pi_0\subset\pi$, \unsure{d'où} \illegible{1} $\pi'_0$ pour $\pi'$, l'hom précédent induit bel et bien un isom
\[
\tag{1}
\mathcal T(\pi_0)_i\simeq\mathcal G(\pi'_0)
\]
d'où
\[
\tag{3}
\widehat{\mathcal T}(\pi_0)\xrightarrow{\sim}\widehat{\mathcal G}(\pi'_0)
\]
et par suite, la suite exacte $\deleted{\unsure{\widehat{\pi}_0}}\widehat{\pi}'_0\longrightarrow\widehat{\mathcal G}(\pi'_0)\longrightarrow\widehat{\mathcal T}(\pi'_0)\longrightarrow 1$ donne
\[
\tag{4}
1\longrightarrow\widehat{\pi}'_0\longrightarrow\widehat{\mathcal T}(\pi_0)\longrightarrow\widehat{\mathcal T}(\pi'_0)\longrightarrow 1
\]
et par suite on trouve un homomorphisme injectif
\[
\begin{gathered}\widehat{\pi}'_0\\[-0.5ex]\Vert\\[-0.5ex]\pi'\end{gathered}
\unsure{\hookrightarrow|\leftrightarrow}
\widehat{\mathcal T}(\pi_0)_i
\hookrightarrow
\begin{gathered}\widehat{\widehat{\mathcal T}}(\pi_0)_i\\[-0.5ex]\Vert\\[-0.5ex]\unsure{\widehat{\widehat{\mathcal T}}}(\pi)\end{gathered}
\]
donc un hom
\[
\tag{5}
i_{\pi_0}:\quad\pi'\longrightarrow\widehat{\widehat{\mathcal T}}(\pi)_i
\]
dont \unsure{le|la} composé avec (1) est l'injection canonique $\pi'\hookrightarrow\widehat{\widehat{\mathcal G}}(\pi')$. Remplaçant la discrétification \inserted{$\pi_0$} par une autre, $\pi_1$, on trouve a priori un autre hom $i_{\pi_1}:\pi'\longrightarrow\widehat{\widehat{\mathcal T}}(\pi)_i$. \unsure{OPS} $\pi_1=u(\pi_0)$, $u\in\widehat{\widehat{\mathcal T}}(\pi)_i$;

\mspage{366}

\scanpage{368}

et par transport de structure on trouve
\[
i_{\pi_1}=i_{u(\pi_0)}=\operatorname{int}(u)\circ i_{\pi_0}\circ u_{\pi'}^{-1}
\]
où $u_{\pi'}$ est l'automorphisme de $\pi'$ défini par $u$ via (1). Dire que $i_{\pi_0}$ est indépendant du choix de la discrétification $\pi_0$, revient aussi à dire que son image dans $\widehat{\widehat{\mathcal T}}(\pi)_i$ est un sous-groupe invariant --- et alors l'action de $\widehat{\widehat{\mathcal T}}(\pi)_i$ sur \deleted{\unsure{$\pi$}} le sous-groupe invariant $i_{\pi_0}(\pi)=i(\pi)$ \deleted{\illegible{1}} \unsure{via} \unsure{autom} intérieurs, n'est autre que celle définie par (1). \uline{Dans} ce cas, on trouve par passage au quotient un \unsure{homomorphisme}
\[
\tag{6}
\widehat{\widehat{\mathcal T}}(\pi)_i/\pi'\longrightarrow\widehat{\widehat{\mathcal T}}(\pi')
\]
dont l'injectivité resp. surjectivité équivaudrait à celle de (1). Mais il n'est pas évident du tout qu'il en soit \unsure{tjrs} ainsi.

S'il n'en était pas ainsi, il \deleted{serait} s'imposerait de regarder le sous-groupe de $\widehat{\widehat{\mathcal T}}(\pi)_i$ \inserted{formé des $\gamma\in\unsure{\widehat{\widehat{\mathcal T}}}(\pi)_i$ qui} \deleted{qui} normalisent le sous-groupe \deleted{\unsure{injecteur|injecté}} $i_{\pi_0}(\pi')$, et tels que l'action induite sur $i_{\pi_0}(\pi')$ correspond à celle donnée par l'action (1) de $\widehat{\widehat{\mathcal T}}(\pi)$ sur $\pi'$. Ce sous-groupe \inserted{$H_{\pi_0,\unsure{i}}$ (qui contient $\widehat{\mathcal T}(\pi_0)_i$) \deleted{\illegible{1}}} dépend donc a priori de la discrétification choisie. Remplaçant $\pi_0$ par $u(\pi_0)$ (où \deleted{\unsure{$u$}} $u\in\widehat{\widehat{\mathcal T}}(\unsure{\pi_0|\pi})_i$) le remplace par $\operatorname{int}(u).H$ --- donc $H$

\mspage{367}

\scanpage{369}
\mspage{187}

tout au moins ne change pas, si on fait varier $\pi_0$ dans \unsure{leur|une} classe de prédiscrétification.

On peut faire des choses plus symétriques, en considérant \uline{pour tout} $i\in I$ le quotient \unsure{correspondant} $\pi'_i$ de $\pi$, d'où, pour une discrétification choisie $\pi_0$, des homomorphismes
\[
\unsure{i_{\pi_0,i}}:\ \pi'_i\longrightarrow
\widehat{\widehat{\mathcal T}}(\pi)_i\subset
\widehat{\widehat{\mathcal T}}(\pi)
\]
et on définira $H_{\pi_0}\subset\widehat{\widehat{\mathcal T}}(\pi)$ comme le sous-groupe des $\gamma\in\widehat{\widehat{\mathcal T}}(\pi)$ qui ``permutent les \deleted{\illegible{1}}\inserted{$i_{\pi_0,i}$} entre eux'' dans un sens évident. On a donc
\[
\tag{7}
H_{\pi_0}\cap\widehat{\widehat{\mathcal T}}(\pi)_i=H_{\pi_0,i}
\]
et
\[
\tag{8}
H\supset\widehat{\mathcal T}(\pi_0)
\]
Le point que j'ai en vue, c'est que le sous-groupe de $\widehat{\widehat{\mathcal T}}(\pi_0)$ image de $\pi_1(M_{g,\nu})$, doit non seulement normaliser $\widehat{\mathcal T}(\pi_0)^+$, mais de plus être contenu dans ce $H$. C'est là une condition que j'ai \unsure{rencontrée} par la bande, \deleted{\unsure{sous}}\inserted{à la faveur de} l'hypothèse (\unsure{peut} être bien hâtive) que (1) est un iso, qui impliquait (tacitement) que l'on avait $H=\widehat{\widehat{\mathcal T}}(\pi_0)$ tout entier. Il est possible que $\widehat{\widehat{\mathcal T}}(\pi_0)$ soit un groupe à tel point \unsure{démesuré|déraisonné} et pathologique, qu'il ne \unsure{pourra|puisse} jamais être question d'en dire des choses raisonnables \inserted{(et vraies)} sur \deleted{\unsure{l'image de}} ce groupe tout entier (\unsure{tel que|telle que} la bijectivité de \unsure{(1)|(4)}, p.\ ex.

\mspage{368}

\scanpage{370}

et qu'on soit obligé de travailler avec des ss-groupes plus petits, qui \unsure{restent} proches du discret (en \unsure{ayant} même des aspects supplémentaires ``arithmétiques'', dûs au $\Gamma_0=\mathrm{Gal}(\unsure{\overline{\mathbb Q}_0}/\mathbb Q)$ !). En fait, il y a \inserted{(pour $I=I(\pi)\ne\varnothing$ i.e.\ $\nu\ne0$)} dans le groupe $\widehat{\mathcal T}(\pi_0)$ une structure \inserted{\unsure{simpliciale}} d'extensions successives, qui va être respectée par l'action extérieure du groupe de Galois, et dont il faudrait tenir compte. Elle fait partie de la ``structure \unsure{à|:} l'$\infty$'' dans les $\pi_1$ des \deleted{\unsure{espaces}} multiplicités modulaires $M_{g,\nu}$, qui même pour $\unsure{\nu|\alpha}=0$ est sans doute non triviale, et il est possible qu'il faille en tenir compte, pour arriver à mettre le doigt sur $\Gamma_{\mathbb Q}$.

Sans essayer de donner d'emblée une description a priori de $\Gamma_{\mathbb Q}$ ``dans les $\widehat{\widehat{\mathcal T}}_{g,\nu}$'', on va procéder de façon plus inductive, en partant de la présence de $\Gamma_{\mathbb Q}$ (pour des raisons \unsure{arithmético-géométriques}), et en essayant de dégager des propriétés de cette présence, peut-être assez fortes pour finir par donner une caractérisation purement algébrique.

\mspage{369}
